% ============================================================================
% Rapport technique — Cloud Optimization (ESILV A5 CCC)
% Compilation : pdfLaTeX (Overleaf : Menu > Compiler = pdfLaTeX)
% Toutes les figures sont dessinées en LaTeX (TikZ / pgfplots) : aucune image
% externe. Les chiffres viennent de docs/RESULTATS.md et du dossier results/.
% ============================================================================
\documentclass[11pt,a4paper]{article}

\usepackage[T1]{fontenc}
\usepackage[utf8]{inputenc}
\usepackage[french]{babel}
\usepackage{lmodern}
\usepackage[a4paper,margin=2.3cm]{geometry}
\usepackage{amsmath,amssymb}
\usepackage{booktabs}
\usepackage{tabularx}
\usepackage{array}
\usepackage{multirow}
\usepackage{xcolor}
\usepackage{float}
\usepackage{caption}
\usepackage{enumitem}
\usepackage{listings}
\usepackage{algorithm}
\usepackage{algpseudocode}
\usepackage{tikz}
\usetikzlibrary{arrows.meta,positioning,shapes.geometric,fit,babel}
\usepackage{pgfplots}
\pgfplotsset{compat=1.17}
\usepackage[hidelinks]{hyperref}

% --- Apparence -------------------------------------------------------------
\definecolor{rrcol}{RGB}{140,140,140}
\definecolor{selcol}{RGB}{31,94,140}
\definecolor{accent}{RGB}{190,80,60}
\captionsetup{font=small,labelfont=bf}
\setlist{itemsep=2pt,topsep=4pt}
\setlength{\parskip}{4pt}
\lstset{basicstyle=\ttfamily\small,columns=fullflexible,breaklines=true,
        frame=single,framerule=0.3pt,rulecolor=\color{gray!50},
        backgroundcolor=\color{gray!6},xleftmargin=4pt,xrightmargin=4pt,
        literate={é}{{\'e}}1 {è}{{\`e}}1 {à}{{\`a}}1 {ê}{{\^e}}1 {ç}{{\c{c}}}1}
\floatname{algorithm}{Algorithme}
\renewcommand{\algorithmicrequire}{\textbf{Entrées :}}
\renewcommand{\algorithmicensure}{\textbf{Sortie :}}

% --- Raccourcis ------------------------------------------------------------
\newcommand{\code}[1]{\texttt{#1}}
\newcommand{\sel}{\texttt{select-maxmin}}
\newcommand{\Drr}{D_{\mathrm{RR}}}
\newcommand{\Crr}{C_{\mathrm{RR}}}
\newcommand{\mesure}{\textsc{mesure}}
\newcommand{\simu}{\textsc{simulation}}

\begin{document}

% ============================================================================
\begin{titlepage}
  \centering
  \vspace*{2cm}
  {\large ESILV — A5 Cloud Computing \& Cybersecurity\par}
  \vspace{1.5cm}
  {\LARGE\bfseries Cloud Optimization\par}
  \vspace{0.4cm}
  {\Large Ordonnancement de tâches distribuées sur Google Cloud\par}
  \vspace{0.3cm}
  {\large Rapport technique du projet semestriel\par}
  \vspace{2cm}
  {\large Prénom NOM \quad·\quad Prénom NOM \quad·\quad Prénom NOM\par}
  \vspace{1cm}
  {\large Novembre 2026\par}
  \vfill
  \begin{minipage}{0.85\textwidth}
    \small
    \textbf{Résumé.} Nous avons construit une petite plateforme qui exécute des
    programmes Go compilés sur quatre machines virtuelles Google Cloud, et un
    \emph{scheduler} qui décide sur quelle machine exécuter chaque tâche. Des
    mesures réelles ont montré que les quatre VM ont la même vitesse pour une
    tâche seule, qu'une VM à 2~vCPU ne vaut qu'environ 1,35~cœur, et que le
    principal levier de coût est le \emph{choix des machines à allumer}. Notre
    algorithme, \sel{}, essaie chaque combinaison « nombre de VM par type » et
    répartit les tâches avec Max-Min. Exécuté réellement sur GCP, il réduit le
    coût de 71\,\% (9~tâches) et de 29\,\% (500~tâches) par rapport à Round
    Robin, et le temps de 24\,\% sur 2000~tâches. Notre modèle de prédiction est
    précis à environ $\pm 10\,\%$ quand les VM fonctionnent normalement ; sa
    principale limite est la variabilité des VM à cœur partagé (e2-medium).
  \end{minipage}
  \vspace{1cm}
\end{titlepage}

\tableofcontents
\newpage

% ============================================================================
\section{Problème et objectifs}
% ============================================================================

\subsection{La question du projet et notre scénario}

Le projet consiste à construire une petite plateforme d'exécution distribuée et
à déterminer, pour un lot de tâches donné, \textbf{sur quelle machine virtuelle
exécuter chaque tâche}.

L'énoncé sépare quatre éléments, que l'on retrouve dans notre architecture
(section~\ref{sec:architecture}) : les \textbf{workloads} (programmes compilés),
l'\textbf{infrastructure} (VM Google Compute Engine), le \textbf{scheduler}
(qui \emph{décide} où va chaque tâche) et l'\textbf{orchestrateur} (qui
\emph{exécute} réellement les tâches et collecte les résultats).

Pour donner un sens à l'optimisation, nous nous plaçons dans la situation
(fictive) d'une \textbf{petite entreprise de traitement par lots}, qui exécute
des calculs pour ses clients : calcul numérique, traitement mémoire, calcul
matriciel, tri de données. Ce n'est pas un service web : un résultat livré
quelques secondes plus tard est acceptable, mais chaque euro dépensé inutilement
se répète à chaque lot. Ses priorités, par ordre d'importance :
\begin{enumerate}
  \item \textbf{réduire la facture cloud} ;
  \item \textbf{conserver un délai raisonnable} ;
  \item pouvoir traiter des lots de quelques tâches à plusieurs milliers ;
  \item rester simple à maintenir et à expliquer.
\end{enumerate}

\subsection{Objectif et contrainte}

Nous évaluons une allocation par deux grandeurs :
\begin{itemize}
  \item $D$, le \textbf{makespan} : le temps entre le début du lot et la fin de la
        dernière tâche ;
  \item $C$, le \textbf{coût} : pour chaque VM utilisée, son temps d'utilisation
        (avec la facturation minimale de 60~s de GCP) multiplié par son prix horaire.
\end{itemize}

Ces deux grandeurs n'ont pas la même unité. Nous les normalisons par celles de
\textbf{Round Robin} (RR), notre référence, et les combinons en un score :
\begin{equation}
  S = 0{,}6 \cdot \frac{C}{\Crr} + 0{,}4 \cdot \frac{D}{\Drr}
  \qquad \text{(plus bas = mieux ; RR vaut 1)}
  \label{eq:score}
\end{equation}

Le poids de 60\,\% sur le coût traduit la priorité de l'entreprise. Ce n'est pas
une valeur imposée : c'est un \textbf{paramètre} du scheduler (option
\code{-wcost}). Pour éviter qu'une solution très économique mais très lente soit
retenue, nous ajoutons une \textbf{contrainte de délai} :
\begin{equation}
  D \le 1{,}25 \cdot \Drr
  \label{eq:deadline}
\end{equation}
Une allocation qui la viole est dite \emph{non admissible}, quel que soit son
score. Le seuil de 125\,\% est un choix de conception, lui aussi paramétrable
(\code{-deadline}).

\subsection{Ce qui a changé depuis notre première version}

Une première version du projet proposait un algorithme « V8 » (glouton pondéré
puis recherche locale), évalué uniquement en simulation, et dont le modèle de
durée sous-estimait les temps réels d'un facteur 3. Nous avons repris le projet
en partant des mesures. Le scénario, le score $S$ et la contrainte de délai sont
conservés ; le reste a été refait :

\begin{table}[H]
\centering\small
\begin{tabularx}{\textwidth}{>{\raggedright\arraybackslash}p{3.3cm}XX}
\toprule
 & \textbf{Première version} & \textbf{Version actuelle} \\
\midrule
Mesures & 180 mesures, VM3 recopiée sur VM2 & 4 VM mesurées séparément + mesure du parallélisme réel \\
Modèle de durée & médiane du benchmark (erreur $\times 3$) & médiane $\times$ contention + surcoût SSH (erreur $\approx \pm 10\,\%$) \\
Algorithme & V8 : recherche locale qui minimise le makespan & \sel{} : choix exact des VM à allumer + Max-Min, minimise $S$ \\
Exécution distribuée & non démontrée & orchestrateur Go, des milliers de tâches exécutées \\
Résultats & simulés uniquement & simulés \emph{et} mesurés sur GCP \\
\bottomrule
\end{tabularx}
\caption{Évolution du projet.}
\end{table}

% ============================================================================
\section{Architecture}
\label{sec:architecture}
% ============================================================================

\subsection{Vue d'ensemble}

La figure~\ref{fig:archi} reprend le flux de données de l'énoncé. La
\textbf{décision} (scheduler) et l'\textbf{exécution} (orchestrateur) sont deux
programmes séparés, reliés par un simple fichier CSV : n'importe quelle
allocation (Round Robin, manuelle, optimisée) peut être exécutée par le même
orchestrateur.

\begin{figure}[H]
\centering
\begin{tikzpicture}[
  node distance=0.55cm and 0.6cm,
  box/.style={draw, rounded corners, align=center, minimum height=0.9cm, font=\small, fill=white},
  file/.style={draw, align=center, font=\small\ttfamily, fill=gray!10, minimum height=0.7cm},
  vm/.style={draw, rounded corners, align=center, font=\small, fill=selcol!12, minimum width=2.2cm, minimum height=0.9cm},
  arr/.style={-{Stealth[length=2.5mm]}, thick}]
  \node[file] (tasks) {tasks.csv};
  \node[file, right=of tasks] (mach) {machines.csv};
  \node[file, right=of mach] (perf) {vm\_performance.csv};
  \node[box, below=0.8cm of mach, fill=selcol!20, minimum width=5cm] (sched) {\textbf{Scheduler} (Go)\\ RR, MET, MCT, Min-Min, Max-Min, \sel};
  \node[file, below=of sched] (alloc) {allocation.csv\\ \normalfont\footnotesize task\_id, workload, size, vm};
  \node[box, below=of alloc, fill=accent!15, minimum width=5cm] (orch) {\textbf{Orchestrateur} (Go)\\ 1 file par VM, N workers, SSH};
  \node[vm, below left=0.9cm and 2.2cm of orch] (v1) {vm1\\ \footnotesize e2-standard-2};
  \node[vm, right=0.35cm of v1] (v2) {vm2\\ \footnotesize e2-medium};
  \node[vm, right=0.35cm of v2] (v3) {vm3\\ \footnotesize e2-medium};
  \node[vm, right=0.35cm of v3] (v4) {highcpu\\ \footnotesize e2-highcpu-4};
  \node[file, below=2.3cm of orch] (res) {results.csv \quad \normalfont\footnotesize(une ligne par tâche, horodatée)};
  \draw[arr] (tasks) -- (sched);
  \draw[arr] (mach) -- (sched);
  \draw[arr] (perf) -- (sched);
  \draw[arr] (sched) -- (alloc);
  \draw[arr] (alloc) -- (orch);
  \foreach \vmnode in {v1,v2,v3,v4} {
    \draw[arr] (orch) -- (\vmnode.north);
    \draw[arr, gray] (\vmnode.south) -- (res);
  }
  \node[font=\footnotesize, gray, right=0.2cm of perf] {benchmarks};
\end{tikzpicture}
\caption{Architecture et flux de données. Les flèches grises représentent la
collecte des résultats par la machine de contrôle.}
\label{fig:archi}
\end{figure}

\subsection{Composants et choix de conception}

Les quatre composants (workloads, infrastructure, scheduler, orchestrateur) et
le langage Go sont fixés par l'énoncé. Le tableau~\ref{tab:composants} indique
où se trouve chacun dans le dépôt.

\begin{table}[H]
\centering\small
\begin{tabularx}{\textwidth}{>{\raggedright\arraybackslash}p{2.6cm}X>{\raggedright\arraybackslash}p{3.4cm}}
\toprule
\textbf{Composant} & \textbf{Rôle} & \textbf{Dans le dépôt} \\
\midrule
Workloads & 4 programmes Go compilés & \code{workloads/} \\
Infrastructure & 4 VM Compute Engine : création, démarrage, arrêt, déploiement, benchmark & \code{scripts/} \\
Scheduler & modèle de prédiction et algorithmes ; écrit \code{allocation.csv} & \code{scheduler/} \\
Orchestrateur & exécute une allocation sur les VM et collecte les résultats & \code{orchestrator/} \\
\bottomrule
\end{tabularx}
\caption{Composants du projet.}
\label{tab:composants}
\end{table}

Dans ce cadre, nous avons fait quatre choix qui nous sont propres :
\begin{itemize}
  \item \textbf{Scheduler et orchestrateur reliés par un simple CSV} : chaque partie
        se teste seule, le scheduler tourne sans GCP, et n'importe quelle allocation
        (Round Robin, manuelle, optimisée) est exécutable par le même orchestrateur.
  \item \textbf{Compilation unique en \code{linux/amd64}} : le poste de contrôle est
        en ARM et les VM en x86 ; on copie des binaires statiques, rien n'est à
        installer sur les VM.
  \item \textbf{SSH avec une connexion réutilisée par VM}, plutôt qu'un outil
        d'orchestration complet : simple à expliquer ; son surcoût par tâche est
        mesuré et réduit (section~\ref{sec:orchestration}).
  \item \textbf{Lots générés avec une graine fixe} : toute l'équipe obtient les mêmes
        lots, donc les mêmes résultats.
\end{itemize}

% ============================================================================
\section{Workloads, infrastructure et mesures}
\label{sec:mesures}
% ============================================================================

Avant d'optimiser, il faut savoir combien de temps prend chaque tâche sur chaque
machine. Toutes les valeurs de cette section sont des \mesure{}s.

\subsection{Workloads}

\begin{table}[H]
\centering\small
\begin{tabular}{llll}
\toprule
\textbf{Workload} & \textbf{Profil} & \textbf{Signification de \code{-{}-size}} & \textbf{Tailles mesurées} \\
\midrule
\code{cpu\_task}    & CPU          & plus grand entier testé (nombres premiers) & 200\,000 ; 500\,000 ; 1\,000\,000 \\
\code{memory\_task} & mémoire      & mémoire allouée et parcourue (Mo)          & 128 ; 256 ; 512 \\
\code{matrix\_task} & CPU + mémoire& dimension $N$ des matrices $N\times N$      & 200 ; 400 ; 600 \\
\code{sort\_task}   & mixte        & nombre d'entiers triés                    & 1\,M ; 5\,M ; 10\,M \\
\bottomrule
\end{tabular}
\caption{Les quatre workloads de référence et leurs trois tailles (celles de l'énoncé).}
\end{table}

Chaque programme affiche une ligne JSON contenant son propre temps de calcul
(\code{runtime\_s}), ce qui permet plus tard de séparer le calcul du surcoût
d'orchestration. Les tailles mesurées sont exactement celles utilisées dans les
lots de tâches : le scheduler lit directement une durée mesurée, sans extrapoler.

\subsection{Infrastructure}

\begin{table}[H]
\centering\small
\begin{tabular}{lllll}
\toprule
\textbf{VM} & \textbf{Type GCP} & \textbf{vCPU} & \textbf{RAM} & \textbf{Prix (\$/h)} \\
\midrule
cloudopt-vm1     & e2-standard-2 & 2            & 8 Go & 0,0737 \\
cloudopt-vm2     & e2-medium     & 2 (partagés) & 4 Go & 0,0369 \\
cloudopt-vm3     & e2-medium     & 2 (partagés) & 4 Go & 0,0369 \\
cloudopt-highcpu & e2-highcpu-4  & 4            & 4 Go & 0,1088 \\
\bottomrule
\end{tabular}
\caption{Les quatre VM (Ubuntu 24.04, x86\_64, zone europe-west1-b). Facturation GCP : à la seconde, minimum 60~s par démarrage.}
\label{tab:vms}
\end{table}

\subsection{Benchmark : une tâche seule}

Protocole (\code{scripts/benchmark.sh}, identique en local et sur chaque VM) :
4 workloads $\times$ 3 tailles $\times$ 5 répétitions, une exécution de chauffe
non mesurée avant chaque taille. Nous retenons la \textbf{médiane}, plus robuste
qu'une moyenne à une exécution perturbée. L'écart-type est inférieur à 1\,\%
dans presque tous les cas.

\begin{table}[H]
\centering\small
\begin{tabular}{llrrrrr}
\toprule
\textbf{Workload} & \textbf{Taille} & \textbf{Local} & \textbf{vm1} & \textbf{vm2} & \textbf{vm3} & \textbf{highcpu} \\
\midrule
cpu\_task    & 200\,000   & 0,003 & 0,045 & 0,047 & 0,045 & 0,045 \\
             & 500\,000   & 0,008 & 0,148 & 0,150 & 0,147 & 0,148 \\
             & 1\,000\,000& 0,019 & 0,378 & 0,383 & 0,378 & 0,379 \\
matrix\_task & 200        & 0,004 & 0,018 & 0,020 & 0,019 & 0,018 \\
             & 400        & 0,023 & 0,101 & 0,105 & 0,101 & 0,104 \\
             & 600        & 0,071 & 0,320 & 0,328 & 0,317 & 0,318 \\
memory\_task & 128 Mo     & 0,039 & 0,206 & 0,224 & 0,207 & 0,222 \\
             & 256 Mo     & 0,078 & 0,406 & 0,435 & 0,402 & 0,427 \\
             & 512 Mo     & 0,149 & 0,797 & 0,854 & 0,793 & 0,846 \\
sort\_task   & 1\,M       & 0,055 & 0,137 & 0,140 & 0,138 & 0,138 \\
             & 5\,M       & 0,294 & 0,720 & 0,733 & 0,719 & 0,726 \\
             & 10\,M      & 0,612 & 1,498 & 1,515 & 1,496 & 1,507 \\
\bottomrule
\end{tabular}
\caption{Médiane du temps d'une tâche seule, en secondes (\code{data/vm\_performance.csv}). Local = VM Linux OrbStack sur Mac Apple Silicon.}
\label{tab:bench}
\end{table}

\paragraph{Constats.}
\begin{enumerate}
  \item \textbf{Les quatre VM ont la même vitesse pour une tâche seule}, à 5\,\%
        près. Un workload n'utilise qu'un cœur, et toutes les VM e2 tournent sur la
        même famille de processeurs. HighCPU n'a pas des cœurs plus rapides : elle
        en a plus. Cela contredit notre première version, où HighCPU semblait
        dominer les autres VM tâche par tâche.
  \item La machine locale est 2,5 à 20 fois plus rapide : elle sert à caractériser
        les workloads, pas à prédire les durées sur GCP.
  \item Sur les VM, une tâche dure de 0,02 à 1,5~s : le surcoût d'orchestration
        n'est pas négligeable.
\end{enumerate}

\subsection{Parallélisme réel : l'expérience « workers »}
\label{sec:workers}

Le benchmark mesure une tâche \emph{seule}. Or l'orchestrateur lance plusieurs
tâches en même temps sur une VM. Nous avons donc mesuré le parallélisme réel
(\code{scripts/exp\_workers.sh}) : le même lot de 24 tâches (2 $\times$ les 12
combinaisons) sur chaque VM, avec 1, 2 puis 4 tâches simultanées
(\emph{workers}), 3 répétitions par niveau.

\begin{figure}[H]
\centering
\begin{tikzpicture}
\begin{axis}[
  width=0.62\textwidth, height=6cm,
  xlabel={Workers (tâches simultanées par VM)}, ylabel={Débit (tâches/s)},
  xtick={1,2,4}, ymin=0, ymax=6, grid=major, grid style={gray!20},
  legend style={font=\small, at={(0.5,-0.2)}, anchor=north, legend columns=2,
                /tikz/every even column/.append style={column sep=0.6cm}}]
  \addplot[mark=*, thick, selcol] coordinates {(1,2.22) (2,3.98) (4,5.32)};
  \addlegendentry{highcpu (4 vCPU)}
  \addplot[mark=square*, thick, accent] coordinates {(1,1.96) (2,2.66) (4,2.72)};
  \addlegendentry{vm1 (2 vCPU)}
  \addplot[mark=triangle*, thick, rrcol] coordinates {(1,2.19) (2,2.97) (4,3.08)};
  \addlegendentry{vm2 (2 vCPU partagés)}
  \addplot[mark=diamond*, thick, gray!50!black, dashed] coordinates {(1,2.22) (2,2.99) (4,2.86)};
  \addlegendentry{vm3 (2 vCPU partagés)}
\end{axis}
\end{tikzpicture}
\caption{Débit mesuré selon le nombre de tâches simultanées (moyenne de 3 répétitions, écart < 3\,\%).}
\label{fig:workers}
\end{figure}

\paragraph{Constats.}
\begin{enumerate}
  \item \textbf{Une VM à 2~vCPU ne vaut qu'environ 1,35 cœur} : passer de 1 à 2
        workers n'augmente le débit que de 35\,\%, car chaque tâche ralentit
        d'environ $\times 1{,}4$. Avec 4 workers, plus aucun gain. Sur GCP, un vCPU
        est un \emph{hyperthread} : 2~vCPU correspondent à un seul cœur physique.
  \item \textbf{HighCPU (4 vCPU, soit 2 cœurs) passe bien à l'échelle} :
        $\times 1{,}8$ avec 2 workers, $\times 2{,}4$ avec 4.
  \item Nous retenons \textbf{workers = nombre de vCPU} (2, ou 4 pour highcpu) :
        c'est le débit maximal ou quasi maximal partout.
\end{enumerate}

\begin{table}[H]
\centering\small
\begin{tabular}{lrrr}
\toprule
\textbf{VM} & \textbf{Débit (tâches/s)} & \textbf{Prix (\$/h)} & \textbf{Débit par \$/h} \\
\midrule
vm2 / vm3 (e2-medium)     & 2,97 / 2,99 & 0,0369 & \textbf{$\approx$ 81} \\
highcpu (e2-highcpu-4)    & 5,32        & 0,1088 & 49 \\
vm1 (e2-standard-2)       & 2,66        & 0,0737 & 36 \\
\bottomrule
\end{tabular}
\caption{Rapport débit/prix au nombre de workers retenu. Les e2-medium sont les plus économiques, highcpu la plus rapide, vm1 est dominée.}
\label{tab:debitprix}
\end{table}

Ce tableau résume le compromis que doit arbitrer le scheduler : \textbf{vm2/vm3
pour économiser, highcpu pour aller vite, vm1 seulement si la capacité manque}.

% ============================================================================
\section{Implémentation distribuée et orchestration}
\label{sec:orchestration}
% ============================================================================

\subsection{Déploiement}

Les VM sont gérées par des scripts \code{gcloud}. Elles sont décrites à un seul
endroit, \code{data/machines.csv} (nom, type GCP, vCPU, RAM, prix, contention),
lu à la fois par ces scripts et par le scheduler : changer de parc de VM revient à
modifier ce fichier. Seuls le projet et la zone GCP sont dans
\code{scripts/config.sh}.

\begin{lstlisting}
./scripts/manage_vms.sh create   # crée les VM manquantes, corrige le type si besoin
./scripts/manage_vms.sh start
./scripts/deploy.sh              # compile en linux/amd64 et copie les binaires
./scripts/manage_vms.sh stop     # toujours à la fin d'une session
\end{lstlisting}

\code{deploy.sh} compile les quatre workloads \textbf{une seule fois}, sur la
machine de contrôle, pour Linux x86\_64 (\code{GOOS=linux GOARCH=amd64}, binaires
statiques), puis les copie sur chaque VM et vérifie qu'ils s'y exécutent. Aucune
installation n'est nécessaire sur les VM.

\subsection{L'orchestrateur}

L'orchestrateur (\code{orchestrator/main.go}) exécute une allocation telle quelle :
il ne déplace jamais une tâche vers une autre VM que celle choisie par le scheduler.

\begin{itemize}
  \item \textbf{Une file par VM}, dans l'ordre de \code{allocation.csv}. Toutes les
        VM travaillent en parallèle ; sur chaque VM, \textbf{N workers} (par défaut
        N = nombre de vCPU, section~\ref{sec:workers}) vident la file : quand un
        worker finit une tâche, il prend la suivante.
  \item \textbf{Une seule connexion SSH par VM}, ouverte au départ et réutilisée par
        toutes ses tâches (multiplexage OpenSSH \code{ControlMaster}), avec la clé et
        la vérification d'identité d'hôte de \code{gcloud}.
  \item \textbf{Horodatage sur la machine de contrôle} : début et fin de chaque
        tâche sont mesurés avec une seule horloge ; aucune synchronisation entre VM
        n'est nécessaire.
  \item \textbf{Séparation calcul / surcoût} : chaque workload renvoie son propre
        temps de calcul ; l'orchestrateur en déduit
        \code{overhead\_s = durée mesurée $-$ runtime\_s}, le coût du réseau et de SSH.
  \item \textbf{Robustesse} : une tâche en échec n'arrête pas les autres (code de
        retour et erreur enregistrés) ; au démarrage, l'orchestrateur attend que le
        serveur SSH des VM soit prêt (une VM \code{RUNNING} refuse encore les
        connexions pendant 20 à 40~s) ; il relève le modèle de processeur de chaque VM.
\end{itemize}

\paragraph{Sorties.} \code{results/<run>.csv} contient une ligne par tâche
(\code{task\_id, workload, size, vm, worker, start\_s, end\_s, duration\_s,
remote\_runtime\_s, overhead\_s, exit\_code, error}) ; \code{results/<run>\_vms.csv}
résume chaque VM (fenêtre active, temps occupé, surcoût moyen, durée facturée,
coût, processeur).

\subsection{Premières preuves d'exécution distribuée}

Premier lot réel : les 9 tâches du TD3 réparties par Round Robin sur les 4 VM.

\begin{table}[H]
\centering\small
\begin{tabular}{lrrrr}
\toprule
\textbf{VM} & \textbf{Tâches} & \textbf{Fenêtre active (s)} & \textbf{Temps de calcul (s)} & \textbf{Surcoût SSH / tâche} \\
\midrule
highcpu & 2 & 0,54 & 0,58 & 0,077 s \\
vm1     & 3 & 1,31 & 1,68 & 0,085 s \\
vm2     & 2 & 0,83 & 1,45 & 0,070 s \\
vm3     & 2 & 0,68 & 0,94 & 0,099 s \\
\bottomrule
\end{tabular}
\caption{Premier lot réel (\mesure{}) : makespan 1,316~s, 9 tâches sur 9 réussies.}
\end{table}

Depuis, l'orchestrateur a exécuté près de 40\,000 tâches sur GCP au cours des
expériences (sections~\ref{sec:workers} et~\ref{sec:reels}), sans perte de tâche.

\paragraph{Le surcoût SSH.} Grâce à la connexion réutilisée, il est de
\textbf{0,06 à 0,12~s par tâche}, contre environ 0,5 à 1~s avec une connexion par
tâche (ordre de grandeur de notre première version). Il augmente avec le nombre
de tâches simultanées (jusqu'à 0,19~s avec 4 workers sur une VM à 2~vCPU) : nous
retenons 0,1~s dans le modèle.

% ============================================================================
\section{Formulation du problème d'ordonnancement}
\label{sec:formulation}
% ============================================================================

\subsection{Données et variables de décision}

\begin{itemize}
  \item $T$ : l'ensemble des tâches du lot ; chaque tâche $i$ a un workload et une taille.
  \item $M$ : l'ensemble des VM ; chaque VM $j$ a un type, $k_j$ workers (= vCPU),
        une RAM et un prix horaire $\pi_j$.
  \item $x_{ij} \in \{0,1\}$ : vaut 1 si la tâche $i$ est exécutée sur la VM $j$.
  \item L'\textbf{ordre} des tâches dans la file de chaque VM, car l'orchestrateur
        les exécute dans l'ordre du fichier. L'ordre fait donc partie de la décision.
\end{itemize}

Une VM qui ne reçoit aucune tâche n'est pas allumée et ne coûte rien : choisir
les $x_{ij}$, c'est aussi choisir \textbf{quelles VM allumer}.

\subsection{Modèle de durée}

La durée prédite de la tâche $i$ sur la VM $j$ combine les trois mesures de la
section~\ref{sec:mesures} :
\begin{equation}
  p_{ij} = \underbrace{\tilde{t}_{ij}}_{\text{médiane du benchmark}}
  \times \underbrace{\gamma_j}_{\text{contention}}
  + \underbrace{\omega}_{\text{surcoût SSH}}
  \label{eq:pij}
\end{equation}
avec $\omega = 0{,}1$~s, $\gamma_j = 1{,}40$ pour les VM à 2~vCPU et $1{,}32$ pour
highcpu.

$\gamma_j$ est mesuré directement dans l'expérience workers
(section~\ref{sec:workers}) : c'est le temps de calcul de chaque tâche, relevé par
le workload lui-même quand tous les workers tournent, divisé par sa médiane de
benchmark (moyenne sur 24 tâches). Il ne se déduit pas du gain de débit
($4/2{,}40 \approx 1{,}67$ pour highcpu), qui inclut aussi le surcoût SSH (non
ralenti par la contention) et les workers inactifs en fin de lot. La valeur est
confirmée en réel : highcpu finit à $+2$ à $+5\,\%$ de sa prédiction dans toutes
nos exécutions (tableau~\ref{tab:precision}).

\subsection{Simulation, makespan et coût}

Pour évaluer une allocation sans l'exécuter, nous simulons exactement ce que fera
l'orchestrateur : chaque VM $j$ dispose de $k_j$ workers ; chaque tâche, dans
l'ordre de la file, est confiée au worker qui se libère le premier. On obtient
l'heure de fin $T_j$ de chaque VM, puis :
\begin{align}
  D &= \max_{j \in M} T_j \\
  C &= \sum_{j \,:\, \text{VM utilisée}} \max(60\ \text{s},\, T_j) \cdot \frac{\pi_j}{3600}
\end{align}
La facturation minimale de 60~s a une conséquence importante : \textbf{pour un
lot de moins d'une minute, le coût ne dépend que des VM allumées}, pas de la
façon dont on répartit les tâches. Au-delà, il dépend aussi du temps passé sur
chaque VM.

\subsection{Contraintes et objectif}

\begin{align}
  & \textstyle\sum_{j \in M} x_{ij} = 1 \quad \forall i \in T
    && \text{(chaque tâche exactement une fois)} \\
  & D \le 1{,}25 \cdot \Drr
    && \text{(délai, éq.~\ref{eq:deadline})} \\
  & k_j \cdot \max_{i : x_{ij}=1} \mathrm{RSS}_{ij} \le \mathrm{RAM}_j \quad \forall j
    && \text{(mémoire)} \\
  & \text{minimiser } S = 0{,}6\, C/\Crr + 0{,}4\, D/\Drr
    && \text{(éq.~\ref{eq:score})}
\end{align}
$\mathrm{RSS}_{ij}$ est la mémoire maximale mesurée de la tâche. Avec nos tailles,
la contrainte mémoire est toujours respectée (4 $\times$ 515~Mo $<$ 4~Go), mais
elle est vérifiée par le programme.

Il faut choisir une VM pour chaque tâche : avec 4 VM et 2000 tâches, cela fait
$4^{2000}$ possibilités (plus de $10^{1200}$), bien trop pour toutes les essayer.
Nous utilisons donc des \textbf{heuristiques} : des méthodes rapides qui
trouvent une bonne solution sans garantir la meilleure. Sur un lot de 9 tâches
($4^9 = 262\,144$ possibilités), on peut en revanche calculer l'optimum exact, ce
qui permet de mesurer l'écart de notre heuristique (section~\ref{sec:simu}).

\subsection{Validation du modèle}

Avant d'utiliser ce modèle pour comparer des algorithmes, nous avons vérifié
qu'il prédit bien la réalité. Pour cela, nous avons exécuté sur GCP un lot de
9 tâches réparti par Round Robin, puis comparé le makespan mesuré à celui prédit
par le modèle.

\begin{table}[H]
\centering\small
\begin{tabular}{lrrr}
\toprule
\textbf{Modèle} & \textbf{Makespan prédit} & \textbf{Makespan mesuré} & \textbf{Écart} \\
\midrule
Médiane du benchmark seule                             & 0,80 s & 1,32 s & réel $\times$ 1,65 \\
Avec contention et surcoût SSH (éq.~\ref{eq:pij})      & 1,22 s & 1,32 s & réel $\times$ 1,08 \\
\bottomrule
\end{tabular}
\caption{Makespan prédit et mesuré sur un lot de 9 tâches (\mesure{}).}
\label{tab:validation}
\end{table}

Avec la médiane seule, le modèle sous-estime le temps réel de 65\,\%. En
ajoutant la contention et le surcoût SSH, l'écart tombe à 8\,\%.

Ce lot n'a pas servi à régler les paramètres $\gamma$ et $\omega$, qui viennent
de l'expérience « workers » : il constitue donc un vrai test. Cette vérification
est intégrée aux tests automatiques du projet
(\code{TestPredictionMatchesRealRun}). La précision du modèle sur les lots de
500 et 2000 tâches est analysée en section~\ref{sec:reels}.

% ============================================================================
\section{Méthodes d'ordonnancement}
\label{sec:algos}
% ============================================================================

Tous les algorithmes utilisent le même modèle (section~\ref{sec:formulation}) et
produisent un \code{allocation.csv} directement exécutable par l'orchestrateur.

\subsection{Algorithmes de référence}

Ce sont les heuristiques classiques de \emph{list scheduling} sur machines
hétérogènes ($n$ tâches, $m$ VM).

\begin{table}[H]
\centering\small
\begin{tabularx}{\textwidth}{>{\raggedright\arraybackslash}p{2.1cm}X>{\raggedright\arraybackslash}p{1.7cm}}
\toprule
\textbf{Algorithme} & \textbf{Principe} & \textbf{Coût} \\
\midrule
Round Robin & tâche $i$ $\rightarrow$ VM $i \bmod m$ ; ignore durées, prix et charge. \textbf{Référence} ($S = 1$). & $O(n)$ \\
MET & chaque tâche va sur la VM où \emph{sa} durée est minimale, sans regarder la file. & $O(nm)$ \\
MCT & dans l'ordre du lot, chaque tâche va sur la VM où elle \emph{finirait} le plus tôt, compte tenu des tâches déjà placées. & $O(nm)$ \\
Min-Min & à chaque tour, on place la tâche dont la meilleure heure de fin est la plus \emph{petite}. & $O(n^2 m)$ \\
Max-Min & à chaque tour, on place la tâche dont la meilleure heure de fin est la plus \emph{grande} (grosses tâches d'abord). & $O(n^2 m)$ \\
\bottomrule
\end{tabularx}
\caption{Heuristiques de référence (\code{scheduler/algorithms/}).}
\end{table}

Ces heuristiques optimisent le \textbf{temps} : aucune ne regarde le prix des VM,
et elles utilisent toujours toutes les machines disponibles. Les simulations de
la section~\ref{sec:simu} montrent que c'est précisément ce qui leur manque.

\subsection{Notre algorithme : \sel}

\paragraph{Idée.} Les mesures ont montré que (1) les VM ont la même vitesse par
tâche, (2) une VM allumée coûte au moins 60~s, et (3) Max-Min répartit déjà bien
les tâches. Le levier principal n'est donc pas la répartition fine, mais
\textbf{le choix des VM à allumer}. De plus, deux VM du même type (vm2 et vm3)
sont interchangeables : on ne se demande pas « quelles VM allumer », mais
\textbf{« combien de VM de chaque type allumer »}.

\paragraph{Méthode, en deux niveaux.}
\begin{enumerate}
  \item \textbf{Choix du nombre de VM par type — exact.} On essaie toutes les
        combinaisons : de 0 à $n_t$ VM pour chaque type $t$, sauf « aucune VM ».
        Pour notre parc : e2-standard-2 (0–1) $\times$ e2-medium (0–2) $\times$
        e2-highcpu-4 (0–1), soit $2 \times 3 \times 2 - 1 = 11$ combinaisons.
  \item \textbf{Répartition — heuristique.} Sur chaque combinaison, Max-Min
        (le meilleur algorithme classique dans nos simulations).
\end{enumerate}
On garde la combinaison de plus petit score $S$ parmi celles qui respectent le
délai et la mémoire.

\begin{algorithm}[H]
\caption{\sel}
\begin{algorithmic}[1]
\Require tâches $T$, VM $M$ regroupées par type, modèle $p_{ij}$
\Ensure une allocation ordonnée
\State $\mathit{ref} \gets$ évaluation de Round Robin sur toutes les VM \Comment{donne $\Drr$, $\Crr$}
\State $\mathit{meilleure} \gets$ aucune
\For{chaque combinaison $(c_1,\dots,c_K)$ avec $0 \le c_t \le n_t$, non toutes nulles}
  \State $V \gets$ les $c_t$ premières VM de chaque type $t$
  \State $A \gets$ Max-Min restreint aux VM de $V$
  \State calculer $D(A)$, $C(A)$ par simulation, puis $S(A)$
  \If{$D(A) \le 1{,}25\,\Drr$ \textbf{et} RAM respectée \textbf{et} $S(A) < S(\mathit{meilleure})$}
    \State $\mathit{meilleure} \gets A$
  \EndIf
\EndFor
\State \Return $\mathit{meilleure}$ (ou Round Robin si aucune combinaison n'est admissible)
\end{algorithmic}
\end{algorithm}

\paragraph{Justification et propriétés.}
\begin{itemize}
  \item \textbf{Décomposition} : choisir les ressources, puis ordonnancer dessus.
        Chaque partie s'explique simplement ; la première est exacte.
  \item \textbf{Jamais pire que Max-Min} (dès que Max-Min est admissible) : la
        combinaison « toutes les VM » fait partie des essais. Si aucune combinaison
        n'est admissible, on renvoie Round Robin, admissible par définition.
        Ces propriétés sont vérifiées par des tests automatiques.
  \item \textbf{Général} : aucun nom de VM, prix ou durée n'est écrit dans le code ;
        tout vient de \code{machines.csv} et \code{vm\_performance.csv}. Le même code
        éteint 2, 1 ou 0 VM selon le lot (section~\ref{sec:simu}). L'objectif reste
        paramétrable (\code{-wcost}, \code{-deadline}).
  \item \textbf{Optimise directement $S$}, contrairement à notre première version,
        dont la recherche locale minimisait le makespan sous un budget de coût.
  \item \textbf{Déterministe et explicable} : \code{./bin/scheduler -explain} affiche
        les 11 combinaisons et leur score (tableau~\ref{tab:explain500}).
\end{itemize}

\paragraph{Complexité et passage à l'échelle.} Le nombre de combinaisons est
$\prod_t (n_t + 1) - 1$ : il dépend du nombre de \emph{types}, pas du nombre de
VM. Avec 3 types et 10 VM de chacun, on obtient 1\,330 combinaisons, contre
environ $10^9$ si l'on traitait chaque VM séparément. Chaque combinaison coûte
un Max-Min en $O(n^2 m)$. Temps mesuré : \textbf{114~ms pour 2000 tâches}. Au-delà
de quelques types, il faudrait remplacer l'énumération par un glouton (ajouter
les types au meilleur débit par dollar jusqu'à tenir le délai).

\paragraph{Intégration.} En mode \code{-algo \sel}, le scheduler affiche aussi
les VM à démarrer (\code{VMS="..." ./scripts/manage\_vms.sh start}) : les VM non
choisies restent éteintes.

% ============================================================================
\section{Résultats en simulation}
\label{sec:simu}
% ============================================================================

\begin{center}
\fbox{\small Tous les chiffres de cette section sont des \textbf{prédictions du modèle} (\simu{}).}
\end{center}

Les lots sont générés par \code{./bin/gentasks -n N -seed 42} : chaque tâche est
tirée uniformément parmi les 12 combinaisons mesurées ; avec la même graine, on
obtient toujours le même fichier. Commande de comparaison :
\code{./bin/scheduler -compare -tasks data/tasks\_N.csv}.

\subsection{Comparaison de tous les algorithmes}

\begin{table}[H]
\centering\small
\begin{tabular}{lrrrcl}
\toprule
\textbf{Algorithme} & \textbf{Coût (\$)} & \textbf{Makespan (s)} & $\boldsymbol{S}$ & \textbf{Admissible} & \textbf{Tâches vm1/vm2/vm3/highcpu} \\
\midrule
\multicolumn{6}{l}{\textit{500 tâches — limite de délai 50,6 s}} \\
Round Robin   & 0,004272 & 40,5  & 1,000 & oui & 125 / 125 / 125 / 125 \\
MET           & 0,003657 & 51,9  & 1,026 & \textbf{non} & 43 / 0 / 92 / 365 \\
MCT           & 0,004272 & 32,5  & 0,921 & oui & 100 / 120 / 104 / 176 \\
Min-Min       & 0,004272 & 32,1  & 0,917 & oui & 100 / 96 / 100 / 204 \\
Max-Min       & 0,004272 & 31,8  & 0,914 & oui & 100 / 96 / 100 / 204 \\
\textbf{\sel} & \textbf{0,003043} & 39,7 & \textbf{0,819} & oui & 0 / 120 / 125 / 255 \\
\midrule
\multicolumn{6}{l}{\textit{2000 tâches — limite de délai 220,5 s}} \\
Round Robin   & 0,009591 & 176,4 & 1,000 & oui & 500 / 500 / 500 / 500 \\
MET           & 0,009976 & 236,2 & 1,160 & \textbf{non} & 183 / 0 / 333 / 1484 \\
MCT           & 0,009701 & 137,1 & 0,918 & oui & 419 / 478 / 423 / 680 \\
Min-Min       & 0,009734 & 136,9 & 0,919 & oui & 398 / 387 / 398 / 817 \\
Max-Min       & 0,009687 & 136,1 & 0,915 & oui & 397 / 388 / 398 / 817 \\
\textbf{\sel} & 0,009687 & 136,1 & \textbf{0,915} & oui & 397 / 388 / 398 / 817 \\
\bottomrule
\end{tabular}
\caption{Comparaison simulée (\code{results/simulation/compare\_*.csv}).}
\label{tab:simu}
\end{table}

\begin{table}[H]
\centering\small
\begin{tabular}{lrrrr}
\toprule
\textbf{Algorithme} & \textbf{Coût (\$)} & \textbf{Makespan (s)} & $\boldsymbol{S}$ & \textbf{VM utilisées} \\
\midrule
Round Robin              & 0,004272 & 1,216 & 1,000 & 4 \\
MET                      & 0,002428 & 1,326 & 0,778 & 2 \\
MCT / Max-Min            & 0,004272 & 1,211 & 0,998 & 4 \\
Min-Min                  & 0,004272 & 1,296 & 1,026 & 4 \\
Allocation manuelle (TD3)& 0,002458 & 1,216 & 0,745 & 3 \\
\textbf{\sel}            & \textbf{0,001230} & 1,266 & \textbf{0,589} & 2 \\
\bottomrule
\end{tabular}
\caption{9 tâches (lot du TD3), avec l'allocation faite à la main au TD3. Limite de délai 1,52 s.}
\label{tab:simu9}
\end{table}

\begin{figure}[H]
\centering
\begin{tikzpicture}
\begin{axis}[
  ybar, width=0.95\textwidth, height=6.2cm, bar width=7pt,
  ylabel={Score $S$ (plus bas = mieux)}, ymin=0.5, ymax=1.2,
  xtick={1,2,3}, xticklabels={9 tâches,500 tâches,2000 tâches},
  xmin=0.5, xmax=3.5, grid=major, grid style={gray!20},
  legend style={at={(0.5,1.02)},anchor=south,legend columns=6,font=\footnotesize},
  extra y ticks={1}, extra y tick style={grid style={black,dashed}}]
  \addplot[fill=rrcol] coordinates {(1,1.000) (2,1.000) (3,1.000)};
  \addplot[fill=accent!60] coordinates {(1,0.778) (2,1.026) (3,1.160)};
  \addplot[fill=gray!30] coordinates {(1,0.998) (2,0.921) (3,0.918)};
  \addplot[fill=gray!55] coordinates {(1,1.026) (2,0.917) (3,0.919)};
  \addplot[fill=gray!80] coordinates {(1,0.998) (2,0.914) (3,0.915)};
  \addplot[fill=selcol] coordinates {(1,0.589) (2,0.819) (3,0.915)};
  \legend{RR, MET, MCT, Min-Min, Max-Min, \sel}
\end{axis}
\end{tikzpicture}
\caption{Score simulé de chaque algorithme. MET n'est pas admissible sur 500 et 2000 tâches.}
\label{fig:simu}
\end{figure}

\paragraph{Constats.}
\begin{enumerate}
  \item \textbf{Les heuristiques classiques optimisent le temps, pas notre
        objectif.} MCT, Min-Min et Max-Min réduisent le makespan d'environ 20\,\%
        en envoyant plus de tâches sur highcpu (4 workers). Max-Min est la meilleure
        des heuristiques classiques ($S \approx 0{,}91$).
  \item \textbf{Mais elles allument toujours les 4 VM} : sous 60~s de lot, leur coût
        reste celui de RR. Sur 9 tâches, l'allocation manuelle du TD3 les bat toutes,
        simplement parce qu'elle n'utilise pas highcpu.
  \item \textbf{MET n'est pas admissible sur les gros lots} : les VM ayant presque
        la même vitesse, choisir la plus rapide tâche par tâche empile tout sur une
        seule machine.
  \item \textbf{\sel{} est le meilleur sur tous les lots.} Par rapport à RR, son
        score baisse de $41\,\%$ sur 9 tâches (2 VM), de $18\,\%$ sur 500 (sans
        vm1) et de $8{,}5\,\%$ sur 2000. Par rapport à Max-Min, la meilleure heuristique
        classique, il gagne $41\,\%$ sur 9 tâches et $10\,\%$ sur 500, et il est
        égal sur 2000 (les 4 VM sont utiles, l'allocation est la même).
\end{enumerate}

\subsection{Pourquoi ces combinaisons ? (\code{-explain})}

\begin{table}[H]
\centering\footnotesize
\begin{tabular}{lrrrc}
\toprule
\textbf{VM allumées} & \textbf{Coût (\$)} & \textbf{Makespan (s)} & $\boldsymbol{S}$ & \textbf{Admissible} \\
\midrule
\textbf{0 × e2-standard-2 + 2 × e2-medium + 1 × e2-highcpu-4} & \textbf{0,003043} & \textbf{39,7} & \textbf{0,819} & \textbf{oui (choisie)} \\
1 × e2-standard-2 + 1 × e2-medium + 1 × e2-highcpu-4 & 0,003657 & 39,7 & 0,906 & oui \\
1 × e2-standard-2 + 2 × e2-medium + 1 × e2-highcpu-4 & 0,004272 & 31,8 & 0,914 & oui \\
0 × e2-standard-2 + 1 × e2-medium + 1 × e2-highcpu-4 & 0,002428 & 52,8 & 0,862 & non \\
1 × e2-standard-2 + 2 × e2-medium + 0 × e2-highcpu-4 & 0,002458 & 53,8 & 0,877 & non \\
0 × e2-standard-2 + 2 × e2-medium + 0 × e2-highcpu-4 & 0,001661 & 81,0 & 1,033 & non \\
\multicolumn{5}{c}{\footnotesize\textit{(5 autres combinaisons, toutes non admissibles)}} \\
\bottomrule
\end{tabular}
\caption{Les combinaisons évaluées par \sel{} sur 500 tâches (limite 50,6 s).}
\label{tab:explain500}
\end{table}

Sur 500 tâches, éteindre vm1 fait gagner 29\,\% de coût pour le même temps que RR.
Deux combinaisons non admissibles ont un meilleur score brut que la combinaison
choisie : \textbf{sans la contrainte de délai, l'algorithme choisirait un lot trop
lent}. La contrainte joue donc réellement son rôle.

Sur 2000 tâches, la combinaison sans vm1 arrive deuxième ($S = 0{,}9235$ contre
$0{,}9146$). Au-delà de 60~s, éteindre vm1 ne supprime pas son travail : les autres
VM tournent plus longtemps. On économise 11\,\% de coût mais le lot dure 25\,\% de
plus ; avec 40\,\% de poids sur le temps, garder vm1 est préférable. Notre intuition
(« éteindre vm1 ») était fausse sur ce lot : c'est l'intérêt de calculer toutes les
combinaisons plutôt que d'appliquer une règle fixe.

% ============================================================================
\section{Résultats réels sur GCP}
\label{sec:reels}
% ============================================================================

\begin{center}
\fbox{\small Tous les chiffres de cette section sont des \textbf{mesures} sur les VM Google Cloud (\mesure{}).}
\end{center}

\subsection{Protocole}

Script \code{scripts/exp\_real.sh}. Pour chaque lot : le scheduler produit
l'allocation de RR et celle de \sel{} (avec leurs prédictions), puis
l'orchestrateur exécute chacune \textbf{3 fois, en alternance} (RR, \sel, RR,
\sel...) : une dérive des VM pendant la session touche ainsi les deux algorithmes.
Les 4 VM sont allumées (RR les utilise toutes) ; \sel{} n'envoie des tâches qu'aux
VM qu'il a choisies. Le coût mesuré applique la même règle que le modèle (fenêtre
active de chaque VM, minimum 60~s). Le score mesuré prend pour référence RR
\emph{mesuré} dans la même session.

Une \textbf{session} correspond à un démarrage des VM suivi de 3 exécutions
alternées par algorithme ; les VM sont arrêtées entre deux sessions. Les sessions
sont numérotées par lot. Comme la vitesse des VM change d'une session à l'autre
(section~\ref{sec:vm3}), les durées ne se comparent qu'à l'intérieur d'une même
session.

\paragraph{Un phénomène à garder en tête.} Pendant nos expériences, une VM
(vm3, une e2-medium) a régulièrement tourné jusqu'à deux fois moins vite que vm2,
de même type et avec la même charge. Nous parlons de VM \textbf{bridée}. Les
exécutions concernées sont signalées dans les tableaux, et la
section~\ref{sec:vm3} explique ce phénomène : c'est une particularité des VM à
cœur partagé de Google Cloud, pas un défaut de notre code.

\subsection{Synthèse : RR contre \sel{}}

\begin{table}[H]
\centering\small
\begin{tabular}{llrrrrrc}
\toprule
\textbf{Lot} & \textbf{Algorithme} & \textbf{Makespan (s)} & \textbf{Coût (\$)} & $\boldsymbol{D/\Drr}$ & $\boldsymbol{S}$ \textbf{mesuré} & $\boldsymbol{S}$ \textbf{prédit} & \textbf{VM} \\
\midrule
\multirow{2}{*}{9 tâches}    & RR    & 1,12 $\pm$ 0,09 & 0,004272 & 1,00 & 1,000 & 1,000 & 4 \\
                             & \sel  & 1,38 $\pm$ 0,15 & \textbf{0,001230} & 1,23 & \textbf{0,664} & 0,589 & 2 \\
\midrule
\multirow{2}{*}{500 tâches}  & RR    & 52,2 $\pm$ 19,2$^{*}$ & 0,004320 & 1,00 & 1,000 & 1,000 & 4 \\
                             & \sel  & 55,8 $\pm$ 16,9$^{*}$ & \textbf{0,003093} & 1,07 & \textbf{0,857} & 0,819 & 3 \\
\midrule
\multirow{2}{*}{2000 tâches} & RR    & 333,6 $\pm$ 6,7 & 0,011819 & 1,00 & 1,000 & 1,000 & 4 \\
                             & \sel  & \textbf{254,3 $\pm$ 2,8} & 0,011453 & 0,76 & \textbf{0,886} & 0,915 & 4 \\
\bottomrule
\end{tabular}
\caption{Moyenne $\pm$ écart-type sur 3 exécutions. Session~1 de chaque lot
(15/11, \code{results/exp\_real/}) ; vm3 était bridée pendant toutes les
exécutions de 2000 tâches. $D/\Drr$ : makespan mesuré divisé par celui de RR
mesuré, admissible si $\le 1{,}25$. $^{*}$ écart-type gonflé par le
ralentissement de vm3 (section~\ref{sec:vm3}).}
\label{tab:reels}
\end{table}

\begin{figure}[H]
\centering
\begin{tikzpicture}
\begin{axis}[
  ybar, width=0.36\textwidth, height=5.5cm, bar width=9pt,
  title={\small Coût mesuré (m\$)}, ymin=0, ymax=14,
  xtick={1,2,3}, xticklabels={9,500,2000}, xlabel={Tâches},
  xmin=0.5, xmax=3.5, grid=major, grid style={gray!20},
  legend style={font=\footnotesize, at={(0.03,0.97)}, anchor=north west}]
  \addplot[fill=rrcol] coordinates {(1,4.272) (2,4.320) (3,11.819)};
  \addplot[fill=selcol] coordinates {(1,1.230) (2,3.093) (3,11.453)};
  \legend{RR, \sel}
\end{axis}
\end{tikzpicture}\hfill
\begin{tikzpicture}
\begin{axis}[
  ybar, width=0.36\textwidth, height=5.5cm, bar width=9pt,
  title={\small Makespan mesuré (RR = 1)}, ymin=0.5, ymax=1.4,
  xtick={1,2,3}, xticklabels={9,500,2000}, xlabel={Tâches},
  xmin=0.5, xmax=3.5, grid=major, grid style={gray!20}]
  \addplot[fill=rrcol] coordinates {(1,1.000) (2,1.000) (3,1.000)};
  \addplot[fill=selcol] coordinates {(1,1.232) (2,1.069) (3,0.762)};
\end{axis}
\end{tikzpicture}\hfill
\begin{tikzpicture}
\begin{axis}[
  ybar, width=0.36\textwidth, height=5.5cm, bar width=9pt,
  title={\small Score $S$ mesuré}, ymin=0.5, ymax=1.1,
  xtick={1,2,3}, xticklabels={9,500,2000}, xlabel={Tâches},
  xmin=0.5, xmax=3.5, grid=major, grid style={gray!20}]
  \addplot[fill=rrcol] coordinates {(1,1.000) (2,1.000) (3,1.000)};
  \addplot[fill=selcol] coordinates {(1,0.664) (2,0.857) (3,0.886)};
\end{axis}
\end{tikzpicture}
\caption{RR contre \sel{} en conditions réelles (session~1) : coût,
makespan relatif à RR et score.}
\label{fig:reels}
\end{figure}

\paragraph{Lecture.}
\begin{itemize}
  \item \textbf{9 tâches} : 2 VM au lieu de 4, \textbf{coût $-71\,\%$}, mais
        lot 23\,\% plus long (modèle : $+4\,\%$), car 9 tâches sur 2 VM à
        2~vCPU ralentissent plus que prévu. La contrainte de délai est respectée
        de justesse ($D/\Drr = 1{,}23 \le 1{,}25$). Score : $-34\,\%$.
  \item \textbf{500 tâches} : vm1 éteinte, \textbf{coût $-29\,\%$}. Les $+7\,\%$
        de temps viennent d'une VM ralentie (section~\ref{sec:vm3}) : sans
        incident, les temps sont égaux (41,4 contre 41,5~s). Score : $-14\,\%$.
  \item \textbf{2000 tâches} : 4 VM mieux réparties, \textbf{temps $-24\,\%$},
        coût $-3\,\%$. Score : $-11\,\%$. Ce gain de temps est en partie dû au
        bridage de vm3 pendant cette session : RR lui confie 500 tâches, \sel{}
        seulement 398. Dans une autre session (tableau~\ref{tab:reels_tous}),
        le gain n'est que de 11\,\%.
\end{itemize}

\subsection{Comparaison réelle de tous les algorithmes}

Pour vérifier que \sel{} bat aussi les heuristiques classiques et pas seulement
Round Robin, nous avons exécuté tous les algorithmes sur GCP (même protocole :
3 exécutions alternées par algorithme ; \code{results/exp\_real\_all/}).

\begin{table}[H]
\centering\small
\begin{tabular}{lrrrrr}
\toprule
\textbf{Algorithme} & \textbf{Makespan mesuré (s)} & \textbf{Coût (\$)} & $\boldsymbol{D/\Drr}$ & $\boldsymbol{S}$ \textbf{mesuré} & $\boldsymbol{S}$ \textbf{prédit} \\
\midrule
\multicolumn{6}{l}{\textit{9 tâches (lot du TD3), session 2}} \\
\textbf{\sel}             & 1,20 $\pm$ 0,05 & \textbf{0,001230} & 0,92 & \textbf{0,541} & 0,589 \\
Allocation manuelle (TD3) & 1,25 $\pm$ 0,13 & 0,002458 & 0,96 & 0,731 & 0,745 \\
MET                       & 1,37 $\pm$ 0,04 & 0,002428 & 1,05 & 0,762 & 0,778 \\
Min-Min                   & 0,87 $\pm$ 0,03 & 0,004272 & 0,67 & 0,868 & 1,026 \\
Max-Min                   & 1,08 $\pm$ 0,07 & 0,004272 & 0,83 & 0,931 & 0,998 \\
MCT                       & 1,10 $\pm$ 0,06 & 0,004272 & 0,84 & 0,937 & 0,998 \\
Round Robin               & 1,30 $\pm$ 0,30 & 0,004272 & 1,00 & 1,000 & 1,000 \\
\midrule
\multicolumn{6}{l}{\textit{500 tâches, session 2}} \\
\textbf{\sel}             & 65,5 $\pm$ 9,3  & \textbf{0,003118} & 1,09 & \textbf{0,864} & 0,819 \\
Min-Min                   & 46,7 $\pm$ 12,2 & 0,004272 & 0,77 & 0,899 & 0,917 \\
Max-Min                   & 47,6 $\pm$ 9,0  & 0,004272 & 0,79 & 0,904 & 0,914 \\
MCT                       & 51,8 $\pm$ 17,0 & 0,004315 & 0,86 & 0,938 & 0,921 \\
Round Robin               & 60,3 $\pm$ 20,1 & 0,004353 & 1,00 & 1,000 & 1,000 \\
MET                       & 73,1 $\pm$ 18,6 & 0,003820 & 1,21 & 1,011 & 1,026 \\
\midrule
\multicolumn{6}{l}{\textit{2000 tâches, session 4$^{\ddagger}$}} \\
\textbf{\sel}             & 162,6 $\pm$ 2,6  & 0,007263 & 0,89 & \textbf{0,969} & 0,915 \\
MCT                       & 165,4 $\pm$ 1,2  & 0,007204 & 0,91 & 0,970 & 0,918 \\
Min-Min                   & 160,6 $\pm$ 3,1  & 0,007360 & 0,88 & 0,973 & 0,919 \\
Max-Min                   & 164,3 $\pm$ 2,1  & 0,007318 & 0,90 & 0,977 & 0,915 \\
Round Robin               & 182,4 $\pm$ 60,4 & \textbf{0,007117} & 1,00 & 1,000 & 1,000 \\
MET                       & 230,3 $\pm$ 14,5 & 0,008019 & \textbf{1,26}$^{\times}$ & 1,181 & 1,160 \\
\bottomrule
\end{tabular}
\caption{Tous les algorithmes exécutés sur GCP (\mesure{}), triés par score
mesuré. Chaque bloc vient d'une seule session, où tous les algorithmes ont été
exécutés en alternance : on compare les algorithmes à l'intérieur d'un bloc,
jamais d'un bloc à l'autre. Sessions du 16/11 (9 et 500 tâches le matin, 2000 tâches le soir),
\code{results/exp\_real\_all/}. $D/\Drr$ : makespan mesuré divisé par celui de
RR mesuré, admissible si $\le 1{,}25$ ; $^{\times}$ non admissible.
$^{\ddagger}$ VM démarrées sur processeur AMD EPYC 7B12, contre Intel Xeon @ 2.20GHz dans toutes les autres sessions (modèle relevé par l'orchestrateur).}
\label{tab:reels_tous}
\end{table}

\paragraph{Comparaison avec le tableau~\ref{tab:reels}.} Les deux tableaux
viennent d'exécutions faites à des moments différents, mais les scores de \sel{}
racontent la même chose : 0,664 puis 0,541 sur 9 tâches, 0,857 puis 0,864 sur
500 tâches, 0,886 puis 0,969 sur 2000 tâches. Il reste toujours sous 1, le score
de RR, et l'écart est le plus grand sur les petits lots, là où éteindre des VM
rapporte le plus.

\begin{figure}[H]
\centering
\begin{tikzpicture}
\begin{axis}[
  ybar, width=0.95\textwidth, height=6cm, bar width=6pt,
  ylabel={Score $S$ mesuré (plus bas = mieux)}, ymin=0.4, ymax=1.25,
  xtick={1,2,3}, xticklabels={9 tâches,500 tâches,2000 tâches}, xmin=0.5, xmax=3.5,
  grid=major, grid style={gray!20},
  legend style={at={(0.5,1.02)}, anchor=south, legend columns=4, font=\footnotesize},
  extra y ticks={1}, extra y tick style={grid style={black,dashed}}]
  \addplot[fill=rrcol]      coordinates {(1,1.000) (2,1.000) (3,1.000)};
  \addplot[fill=accent!60]  coordinates {(1,0.762) (2,1.011) (3,1.181)};
  \addplot[fill=gray!30]    coordinates {(1,0.937) (2,0.938) (3,0.970)};
  \addplot[fill=gray!55]    coordinates {(1,0.868) (2,0.899) (3,0.973)};
  \addplot[fill=gray!80]    coordinates {(1,0.931) (2,0.904) (3,0.977)};
  \addplot[fill=orange!60]  coordinates {(1,0.731) (2,0) (3,0)};
  \addplot[fill=selcol]     coordinates {(1,0.541) (2,0.864) (3,0.969)};
  \legend{RR, MET, MCT, Min-Min, Max-Min, Manuelle, \sel}
\end{axis}
\end{tikzpicture}
\caption{Score mesuré sur GCP. L'allocation manuelle n'existe que pour le lot de 9 tâches.}
\label{fig:reels_tous}
\end{figure}

\begin{figure}[H]
\centering
% Généré par ./bin/gantt à partir de : results/exp_real_all/tasks/rr_r1.csv results/exp_real_all/tasks/select-maxmin_r1.csv
% Durées MESURÉES par l'orchestrateur (start_s / end_s).
\begin{tikzpicture}[x=6.8571cm, y=0.42cm, font=\scriptsize]
  \node[anchor=south west, font=\small\bfseries] at (0,0.40) {Round Robin};
  \draw[gray!25] (0.0000,-11.20) -- (0.0000,-0.20);
  \draw[gray!25] (0.2500,-11.20) -- (0.2500,-0.20);
  \draw[gray!25] (0.5000,-11.20) -- (0.5000,-0.20);
  \draw[gray!25] (0.7500,-11.20) -- (0.7500,-0.20);
  \draw[gray!25] (1.0000,-11.20) -- (1.0000,-0.20);
  \draw[gray!25] (1.2500,-11.20) -- (1.2500,-0.20);
  \draw[gray!25] (1.5000,-11.20) -- (1.5000,-0.20);
  \draw[gray!25] (1.7500,-11.20) -- (1.7500,-0.20);
  \node[anchor=east] at (0,-1.00) {vm1 w1};
  \node[anchor=east] at (0,-2.00) {vm1 w2};
  \node[anchor=east] at (0,-3.25) {vm2 w1};
  \node[anchor=east] at (0,-4.25) {vm2 w2};
  \node[anchor=east] at (0,-5.50) {vm3 w1};
  \node[anchor=east] at (0,-6.50) {vm3 w2};
  \node[anchor=east] at (0,-7.75) {highcpu w1};
  \node[anchor=east] at (0,-8.75) {highcpu w2};
  \node[anchor=east] at (0,-9.75) {highcpu w3};
  \node[anchor=east] at (0,-10.75) {highcpu w4};
  \draw[fill=green!45, draw=black!60] (0.0002,-9.07) rectangle (0.1814,-8.43);
  \node at (0.0908,-8.75) {T4};
  \draw[fill=orange!50, draw=black!60] (0.0005,-10.07) rectangle (0.5011,-9.43);
  \node at (0.2508,-9.75) {T8};
  \draw[fill=blue!35, draw=black!60] (0.0008,-3.57) rectangle (0.6016,-2.93);
  \node at (0.3012,-3.25) {T2};
  \draw[fill=blue!35, draw=black!60] (0.0008,-2.32) rectangle (0.4330,-1.68);
  \node at (0.2169,-2.00) {T1};
  \draw[fill=orange!50, draw=black!60] (0.0008,-1.32) rectangle (1.6443,-0.68);
  \node at (0.8226,-1.00) {T5};
  \draw[fill=orange!50, draw=black!60] (0.0009,-6.82) rectangle (0.6370,-6.18);
  \node at (0.3190,-6.50) {T3};
  \draw[fill=green!45, draw=black!60] (0.0010,-4.57) rectangle (0.6341,-3.93);
  \node at (0.3176,-4.25) {T6};
  \draw[fill=blue!35, draw=black!60] (0.0012,-5.82) rectangle (0.2690,-5.18);
  \node at (0.1351,-5.50) {T7};
  \draw[fill=green!45, draw=black!60] (0.4331,-2.32) rectangle (0.8921,-1.68);
  \node at (0.6626,-2.00) {T9};
  \draw[->] (0,-11.20) -- (1.8333,-11.20) node[right] {s};
  \draw (0.0000,-11.20) -- (0.0000,-11.35) node[below] {0};
  \draw (0.2500,-11.20) -- (0.2500,-11.35) node[below] {0{,}25};
  \draw (0.5000,-11.20) -- (0.5000,-11.35) node[below] {0{,}5};
  \draw (0.7500,-11.20) -- (0.7500,-11.35) node[below] {0{,}75};
  \draw (1.0000,-11.20) -- (1.0000,-11.35) node[below] {1};
  \draw (1.2500,-11.20) -- (1.2500,-11.35) node[below] {1{,}25};
  \draw (1.5000,-11.20) -- (1.5000,-11.35) node[below] {1{,}5};
  \draw (1.7500,-11.20) -- (1.7500,-11.35) node[below] {1{,}75};
  \node[anchor=south west, font=\small\bfseries] at (0,-12.40) {select-maxmin};
  \draw[gray!25] (0.0000,-24.00) -- (0.0000,-13.00);
  \draw[gray!25] (0.2500,-24.00) -- (0.2500,-13.00);
  \draw[gray!25] (0.5000,-24.00) -- (0.5000,-13.00);
  \draw[gray!25] (0.7500,-24.00) -- (0.7500,-13.00);
  \draw[gray!25] (1.0000,-24.00) -- (1.0000,-13.00);
  \draw[gray!25] (1.2500,-24.00) -- (1.2500,-13.00);
  \draw[gray!25] (1.5000,-24.00) -- (1.5000,-13.00);
  \draw[gray!25] (1.7500,-24.00) -- (1.7500,-13.00);
  \node[anchor=east] at (0,-13.80) {vm1 w1};
  \fill[gray!12] (0,-14.18) rectangle (1.7500,-13.42);
  \node[anchor=east] at (0,-14.80) {vm1 w2};
  \fill[gray!12] (0,-15.18) rectangle (1.7500,-14.42);
  \node[gray] at (0.8750,-14.30) {VM éteinte};
  \node[anchor=east] at (0,-16.05) {vm2 w1};
  \node[anchor=east] at (0,-17.05) {vm2 w2};
  \node[anchor=east] at (0,-18.30) {vm3 w1};
  \node[anchor=east] at (0,-19.30) {vm3 w2};
  \node[anchor=east] at (0,-20.55) {highcpu w1};
  \fill[gray!12] (0,-20.93) rectangle (1.7500,-20.17);
  \node[anchor=east] at (0,-21.55) {highcpu w2};
  \fill[gray!12] (0,-21.93) rectangle (1.7500,-21.17);
  \node[anchor=east] at (0,-22.55) {highcpu w3};
  \fill[gray!12] (0,-22.93) rectangle (1.7500,-22.17);
  \node[anchor=east] at (0,-23.55) {highcpu w4};
  \fill[gray!12] (0,-23.93) rectangle (1.7500,-23.17);
  \node[gray] at (0.8750,-22.05) {VM éteinte};
  \draw[fill=orange!50, draw=black!60] (0.0001,-19.62) rectangle (1.1738,-18.98);
  \node at (0.5870,-19.30) {T5};
  \draw[fill=orange!50, draw=black!60] (0.0001,-17.37) rectangle (0.6416,-16.73);
  \node at (0.3209,-17.05) {T8};
  \draw[fill=orange!50, draw=black!60] (0.0001,-18.62) rectangle (0.6646,-17.98);
  \node at (0.3324,-18.30) {T3};
  \draw[fill=blue!35, draw=black!60] (0.0002,-16.37) rectangle (0.5891,-15.73);
  \node at (0.2946,-16.05) {T2};
  \draw[fill=green!45, draw=black!60] (0.5892,-16.37) rectangle (1.1406,-15.73);
  \node at (0.8649,-16.05) {T6};
  \draw[fill=blue!35, draw=black!60] (0.6417,-17.37) rectangle (0.8861,-16.73);
  \node at (0.7639,-17.05) {T7};
  \draw[fill=blue!35, draw=black!60] (0.6647,-18.62) rectangle (0.9160,-17.98);
  \node at (0.7903,-18.30) {T1};
  \draw[fill=green!45, draw=black!60] (0.8862,-17.37) rectangle (1.1150,-16.73);
  \node at (1.0006,-17.05) {T9};
  \draw[fill=green!45, draw=black!60] (0.9160,-18.62) rectangle (1.1348,-17.98);
  \node at (1.0254,-18.30) {T4};
  \draw[->] (0,-24.00) -- (1.8333,-24.00) node[right] {s};
  \draw (0.0000,-24.00) -- (0.0000,-24.15) node[below] {0};
  \draw (0.2500,-24.00) -- (0.2500,-24.15) node[below] {0{,}25};
  \draw (0.5000,-24.00) -- (0.5000,-24.15) node[below] {0{,}5};
  \draw (0.7500,-24.00) -- (0.7500,-24.15) node[below] {0{,}75};
  \draw (1.0000,-24.00) -- (1.0000,-24.15) node[below] {1};
  \draw (1.2500,-24.00) -- (1.2500,-24.15) node[below] {1{,}25};
  \draw (1.5000,-24.00) -- (1.5000,-24.15) node[below] {1{,}5};
  \draw (1.7500,-24.00) -- (1.7500,-24.15) node[below] {1{,}75};
  \draw[fill=blue!35, draw=black!60] (0.0000,-25.60) rectangle ++(0.0583,0.5);
  \node[anchor=west] at (0.0656,-25.35) {cpu\_task};
  \draw[fill=orange!50, draw=black!60] (0.3792,-25.60) rectangle ++(0.0583,0.5);
  \node[anchor=west] at (0.4448,-25.35) {memory\_task};
  \draw[fill=green!45, draw=black!60] (0.7583,-25.60) rectangle ++(0.0583,0.5);
  \node[anchor=west] at (0.8240,-25.35) {matrix\_task};
\end{tikzpicture}

\caption{Diagramme temporel de la première exécution de chaque
algorithme sur 9 tâches : une ligne par worker, une barre par tâche (début et fin
relevés par l'orchestrateur). Round Robin allume les 4 VM et laisse la moitié des
workers de highcpu inactifs ; \sel{} n'allume que vm2 et vm3 et occupe leurs deux
workers en parallèle pendant tout le lot.}
\label{fig:gantt}
\end{figure}

\paragraph{Lecture.}
\begin{itemize}
  \item \textbf{Sur 9 et 500 tâches, \sel{} obtient le meilleur score mesuré},
        devant les cinq heuristiques et l'allocation manuelle : il éteint des VM,
        ce que les autres ne font pas.
  \item Sur 500 tâches, \textbf{cette victoire vient entièrement du coût}
        ($-28\,\%$). Avec 3 exécutions et des écarts-types de 9 à 20~s, les écarts
        de makespan ne sont pas significatifs : son lot est même plus long, car
        vm3 a été bridée pendant 2 de ses 3 exécutions.
  \item \textbf{Sur 2000 tâches, \sel{} est à égalité avec Max-Min, MCT et
        Min-Min} (écarts inférieurs à 1\,\%, dans le bruit de mesure). Il garde
        les 4 VM : son allocation est alors exactement celle de Max-Min.
  \item \textbf{La contrainte de délai est aussi respectée en réel} par \sel{} sur
        tous les lots (au plus 1,23, sur 9 tâches en session 1). Seul MET la
        dépasse, sur 2000 tâches (1,26), ce que le modèle avait prévu. Sur
        500 tâches, le modèle prévoyait aussi MET non admissible (51,9~s pour
        une limite de 50,6~s), mais il passe en réel (1,21) : RR, ralenti par
        vm3 dans cette session, a relâché la contrainte. C'est la limite d'une
        échéance relative à RR (section~\ref{sec:limites}).
  \item \textbf{Le classement mesuré suit globalement le classement prédit},
        sauf Min-Min sur 9 tâches : prédit avec le pire score, il est en réalité
        le plus rapide (section~\ref{sec:limites}).
\end{itemize}

\subsection{Précision du modèle}

\begin{table}[H]
\centering\small
\begin{tabular}{llr@{\hspace{1.2em}}ccc@{\hspace{1.2em}}ccc}
\toprule
 & & \textbf{Prédit} & \multicolumn{3}{c}{\textbf{Mesuré (s)}} & \multicolumn{3}{c}{\textbf{Mesuré / prédit}} \\
\cmidrule(lr){4-6}\cmidrule(lr){7-9}
\textbf{Lot} & \textbf{Session} & \textbf{(s)} & r1 & r2 & r3 & r1 & r2 & r3 \\
\midrule
9 tâches    & 1 & 1,27  & 1,54  & 1,23  & 1,36  & 1,22 & 0,97 & 1,08 \\
            & 2 & 1,27  & 1,17  & 1,16  & 1,26  & 0,93 & 0,91 & 0,99 \\
\midrule
500 tâches  & 1 & 39,7  & 41,5  & 51,4$^{b}$  & 74,5$^{b}$  & \textbf{1,05} & 1,30 & 1,88 \\
            & 2 & 39,7  & 70,8$^{b}$  & 71,0$^{b}$  & 54,8$^{b}$  & 1,78 & 1,79 & 1,38 \\
\midrule
2000 tâches & 1 & 136,1 & 256,8$^{b}$ & 254,6$^{b}$ & 251,3$^{b}$ & 1,89 & 1,87 & 1,85 \\
            & 2 & 136,1 & 255,5$^{b}$ & 253,7$^{b}$ & 256,2$^{b}$ & 1,88 & 1,86 & 1,88 \\
            & 3 & 136,1 & 263,9$^{b}$ & 262,2$^{b}$ & 148,1 & 1,94 & 1,93 & \textbf{1,09} \\
            & 4$^{\ddagger}$ & 136,1 & 161,0$^{b}$ & 165,7$^{b}$ & 161,3$^{b}$ & 1,18 & 1,22 & 1,19 \\
\bottomrule
\end{tabular}
\caption{Makespan de \sel{} : prédit par le modèle et mesuré sur GCP, pour
chacune des 24 exécutions. $^{b}$ vm3 bridée pendant l'exécution (plus de 1,3 fois
plus lente que vm2 pour la même charge). $^{\ddagger}$ VM démarrées sur processeur AMD EPYC 7B12, contre Intel Xeon @ 2.20GHz dans toutes les autres sessions (modèle relevé par l'orchestrateur).
Sources : \code{results/exp\_real/} et \code{results/exp\_real\_all/}
(\code{predicted.csv}, \code{runs.csv}, \code{*\_vms.csv}).}
\label{tab:precision_select}
\end{table}

\paragraph{Lecture.}
\begin{itemize}
  \item \textbf{Quand aucune VM n'est bridée, le modèle est précis} : le rapport
        mesuré/prédit va de 0,91 à 1,22 (9 tâches, 500 tâches r1, 2000 tâches
        session 3 r3), soit une erreur d'environ $\pm 10$ à 20\,\%.
  \item \textbf{Quand vm3 est bridée, le lot dure jusqu'à 1,9 fois plus longtemps
        que prévu} : le makespan est fixé par la VM la plus lente, et le modèle ne
        peut pas prévoir ce bridage.
  \item Sur 9 tâches, chaque VM n'exécute que quelques tâches, en moins de 2~s :
        aucune n'a le temps d'être bridée.
  \item En session~4 des 2000 tâches, le rapport reste à 1,2 malgré le bridage de
        vm3, parce que les VM tournaient sur processeur AMD, plus rapide que les
        Intel sur lesquels le modèle a été calibré : les deux effets se compensent
        en partie.
\end{itemize}

\paragraph{D'où vient l'erreur ?} Pour le savoir, nous avons comparé l'heure de fin
mesurée et prédite de chaque VM séparément (session 1).

\begin{table}[H]
\centering\small
\begin{tabular}{llcccc}
\toprule
\textbf{Lot} & \textbf{Algorithme} & \textbf{highcpu} & \textbf{vm1} & \textbf{vm2} & \textbf{vm3} \\
\midrule
9    & RR    & 0,72 & 0,89 & 1,01 & 0,93 \\
9    & \sel  & —    & —    & 0,92 & 1,12 \\
500  & RR    & 1,02 & 1,09 & 0,92 & 0,99 \\
500  & \sel  & 1,05 & —    & 0,95 & 1,30 \\
2000 & RR    & 1,02 & 1,09 & 1,15 & \textbf{1,91} \\
2000 & \sel  & 1,05 & 1,15 & 0,96 & \textbf{1,87} \\
\bottomrule
\end{tabular}
\caption{Heure de fin mesurée / prédite pour chaque VM (session~1, médiane des
3 exécutions).}
\label{tab:precision}
\end{table}

Sur 500 et 2000 tâches, hors vm3, le modèle se trompe de $-8$ à $+15\,\%$
(highcpu : $+2$ à $+5\,\%$) ; sur 9 tâches, les écarts sont plus grands (2 ou
3 tâches par VM). vm3 a fortement ralenti sur les lots longs : 1,7 à 1,9 fois
plus lente que prévu sur 2000 tâches, quel que soit l'algorithme. Or le makespan
est fixé par la VM la plus lente : \textbf{une seule VM ralentie suffit à fausser
la prédiction du lot entier} (section~\ref{sec:vm3}).

\subsection{Investigation : le ralentissement des e2-medium}
\label{sec:vm3}

Sur les lots longs, le makespan était souvent bien supérieur à la prédiction.
Pour comprendre pourquoi, nous avons regardé le temps de travail de chaque VM
séparément. Avec \sel{}, l'allocation de 2000 tâches est la même à chaque
session : 388 tâches sur vm2 et 398 sur vm3, deux e2-medium. Elles devraient donc
finir presque en même temps.

\begin{table}[H]
\centering\small
\begin{tabular}{llcc}
\toprule
\textbf{Session} & \textbf{vm3} & \textbf{vm2 (s) : r1 / r2 / r3} & \textbf{vm3 (s) : r1 / r2 / r3} \\
\midrule
1 & d'origine & 130 / 130 / 132 & 257 / 255 / 251 \\
2 & d'origine & 165 / 132 / 132 & 256 / 254 / 256 \\
3 & recréée   & 172 / 131 / 131 & 264 / 262 / 148 \\
4$^{\ddagger}$ & recréée &  87 /  92 /  88 & 161 / 166 / 161 \\
\bottomrule
\end{tabular}
\caption{Temps de travail de vm2 et vm3 sur les 12 exécutions de \sel{} sur
2000 tâches (\mesure{}). Même type de VM, même charge à 3\,\% près.
$^{\ddagger}$ VM démarrées sur processeur AMD EPYC 7B12, contre Intel Xeon @ 2.20GHz dans toutes les autres sessions (modèle relevé par l'orchestrateur).}
\label{tab:vm2vm3}
\end{table}

Les sessions 2 et 3 ont été lancées pour étudier vm3, pas pour comparer les
algorithmes : elles ne figurent donc pas dans les tableaux de synthèse. Leurs
exécutions de RR (241 à 348~s) sont dans \code{results/exp\_real/tasks\_2000\_session2/}
et \code{tasks\_2000\_vm3neuve/} ; la dernière de la session 3 est exclue (surcoût
SSH multiplié par 5 sur les 4 VM à la fois, incident côté machine de contrôle).
Sur ces 5 paires valides, \sel{} est plus rapide que RR dans 3. Les 2 autres sont
les premières exécutions de chaque session : vm3 n'était pas encore bridée pendant
l'exécution de RR, mais l'était pendant celle de \sel{}. En moyenne sur chaque
session, \sel{} reste plus rapide : $-17\,\%$ en session~2 et $-11\,\%$ en
session~3 (sur les paires valides).

\paragraph{Lecture.}
\begin{itemize}
  \item \textbf{vm3 met presque deux fois plus de temps que vm2} dans 11 exécutions
        sur 12. Le temps de calcul mesuré par les workloads eux-mêmes double : c'est
        le processeur qui est ralenti, pas le réseau.
  \item \textbf{Ce n'est pas une machine défectueuse} : nous avons supprimé puis
        recréé vm3, qui s'est probablement retrouvée sur un autre serveur
        physique. Le ralentissement est resté le même (sessions 3 et 4).
  \item vm2 n'a connu que deux épisodes légèrement lents (165 et 172~s), tous deux
        en début de session.
\end{itemize}

\paragraph{Explication.} La documentation de
GCP\footnote{\url{https://cloud.google.com/compute/docs/general-purpose-machines\#cpu-bursting}}
indique que chacun des 2~vCPU d'une e2-medium n'a droit qu'à 50\,\% du temps
processeur, soit l'équivalent d'un seul vCPU. Au-delà, la VM fonctionne par
\emph{rafales} limitées dans le temps. Privée de rafale, elle tourne donc environ
deux fois moins vite, ce que nous mesurons. Avec Round Robin, vm3 n'était
d'ailleurs normale qu'à la \emph{première} exécution de chaque session, puis
ralentissait, ce qui va dans le sens d'une réserve de rafale qui s'épuise. Ce
mécanisme n'explique pas pourquoi vm3 est bridée et pas vm2 : cela dépend
probablement de l'occupation du serveur physique, que nous ne contrôlons pas.
highcpu a été stable dans toutes nos mesures. vm1 (sans cœur partagé) a ralenti
lors d'une session (47 à 66~s au lieu de $\approx 31$~s sur 500 tâches).

\paragraph{Conséquence : nous dépendons de Google.} Sur un lot long, le makespan
est fixé par la VM la plus lente, quel que soit l'algorithme. Or la vitesse réelle
d'une VM dépend de deux choses que le client ne voit pas et ne choisit pas :
\begin{itemize}
  \item \textbf{le serveur physique attribué} : GCP choisit lui-même le processeur
        à chaque démarrage d'une VM e2. Nous l'avons constaté : toutes nos sessions
        ont tourné sur Intel Xeon @ 2.20GHz, sauf la session~4 des 2000 tâches, sur
        AMD EPYC 7B12, où les VM non bridées étaient environ 35\,\% plus rapides
        (highcpu : 92~s au lieu de 143~s pour les mêmes 817 tâches).
        L'orchestrateur l'a enregistré car il relève le modèle de processeur à
        chaque exécution ;
  \item \textbf{l'activité des autres clients sur ce serveur} : une e2-medium ne
        dépasse sa part garantie que si le serveur a du processeur libre. Deux VM
        identiques peuvent donc tourner à des vitesses très différentes, comme vm2
        et vm3. Nous ne pouvons pas le prouver directement (le serveur physique
        n'est pas visible), mais c'est l'explication la plus cohérente avec nos
        mesures.
\end{itemize}
Quand une VM est bridée, toute prédiction faite à l'avance est donc fausse d'un
facteur proche de 2. Ce n'est pas un défaut du modèle qu'un meilleur calibrage
corrigerait : c'est une incertitude propre au cloud mutualisé, qu'un scheduler
doit accepter et gérer. C'est la principale limite de notre approche
(section~\ref{sec:limites}).

% ============================================================================
\section{Discussion et limites}
\label{sec:limites}
% ============================================================================

\subsection{Ce que les mesures nous ont appris}
\begin{enumerate}
  \item \textbf{Mesurer avant d'optimiser change le problème.} Deux idées
        intuitives se sont révélées fausses : HighCPU n'est pas plus rapide
        par tâche, et une VM à 2~vCPU n'exécute pas 2 tâches à pleine vitesse.
        En tenir compte a fait passer l'erreur du modèle de 65\,\% à 8\,\%
        (tableau~\ref{tab:validation}), et a montré que le vrai levier est le
        choix des VM à allumer.
  \item \textbf{Le choix des VM compte plus que la répartition des tâches},
        surtout pour les petits lots, à cause de la facturation minimale de
        60~s.
  \item \textbf{Une règle fixe ne suffit pas} : éteindre vm1 est le bon choix
        sur 500 tâches, mais pas sur 2000. Seul le calcul de toutes les
        combinaisons permet de trancher.
\end{enumerate}

\subsection{Limites}

\paragraph{Variabilité des VM à cœur partagé.} C'est la limite principale. Les
e2-medium sont les VM les moins chères, et donc celles que préfère notre
algorithme. Mais GCP ne leur garantit qu'un vCPU sur deux : au-delà, elles
dépendent de rafales limitées, et leur vitesse peut chuter de moitié sous charge
prolongée (section~\ref{sec:vm3}). Notre modèle, calibré sur des mesures courtes
(tableau~\ref{tab:debitprix}), ne le prévoit pas, et un planning calculé à
l'avance ne s'y adapte pas.

Nous avons simulé ce pire cas en modélisant les deux e2-medium à mi-vitesse
($\gamma = 2{,}8$ au lieu de 1,4). \sel{} change alors de choix : highcpu seule
sur 9 tâches, une seule e2-medium et highcpu sur 500 tâches, toujours les 4 VM
sur 2000 tâches. Sur 9 tâches, ce scénario ne se produit pas en pratique : aucune
VM n'a été bridée sur un lot aussi court. Sur les lots longs, un modèle à deux
vitesses (pleine vitesse pendant la rafale, puis vitesse garantie) donnerait des
choix plus prudents.

\paragraph{Planning statique.} Le scheduler décide tout avant l'exécution ; si une
VM ralentit, ses tâches attendent même si d'autres VM ont fini.

\paragraph{Peu de répétitions.} Chaque algorithme n'a été exécuté que 3 fois par
session. C'est suffisant pour les écarts de coût, stables d'une exécution à
l'autre, mais pas pour départager des makespans proches quand une VM est bridée.

\paragraph{Contention appliquée partout.} Le modèle ralentit toutes les tâches
de $\times 1{,}4$, comme si chaque VM était pleinement occupée. C'est juste sur les
gros lots, mais trop pessimiste sur les petits, où une VM peut n'exécuter qu'une
tâche : sur 9 tâches, Min-Min (une seule tâche sur vm2), prédit le plus lent,
s'est révélé le plus rapide.

\paragraph{Contrainte de délai relative.} L'échéance $D \le 1{,}25 \cdot \Drr$
dépend de RR : plus RR est lent, plus elle est lâche. Elle sert à comparer les
algorithmes sur des lots de tailles très différentes ; pour un client réel, on la
remplacerait par son échéance contractuelle, sans changer l'algorithme.

\paragraph{Coût calculé, pas facture réelle.} Notre coût ne compte que le temps
de travail des VM (au moins 60~s), sans le démarrage, les temps morts, le
disque ni l'IP. Il sert à \emph{comparer} les allocations, pas à chiffrer une
économie réelle.

\subsection{Pistes d'amélioration}

\begin{itemize}
  \item \textbf{Modéliser deux vitesses pour les e2-medium} : pleine vitesse
        pendant la rafale, environ moitié au-delà. Les prédictions des lots longs
        seraient plus prudentes et \sel{} enverrait moins de tâches aux e2-medium.
  \item \textbf{Mesurer la vitesse réelle des VM juste avant de planifier} (quelques
        tâches de test), et mettre à jour le modèle.
  \item \textbf{Rééquilibrage dynamique} (\emph{work stealing}) : une VM qui a fini
        prend des tâches restantes d'une VM lente, au prix d'un écart au planning.
  \item \textbf{Surveiller le \emph{steal time}} des VM pour confirmer et quantifier
        le bridage des e2-medium.
  \item \textbf{Recherche locale} après \sel{} (déplacements et échanges qui font
        baisser $S$), pour affiner la répartition au-delà de Max-Min.
  \item \textbf{Plus de répétitions} (10 par algorithme) pour départager les
        makespans proches.
\end{itemize}

% ============================================================================
\section{Conclusion}
% ============================================================================

Nous avons construit une chaîne complète et reproductible : des workloads Go
compilés une fois et déployés sur quatre VM Google Cloud, un orchestrateur qui
exécute des files de tâches en parallèle et mesure chaque tâche, et un scheduler
qui décide où exécuter chaque tâche à partir d'un modèle calibré sur nos mesures.

Les mesures ont guidé l'optimisation. Elles ont montré que les VM ont la même
vitesse par tâche, qu'une VM à 2~vCPU ne vaut qu'environ 1,35 cœur, et que le coût
dépend avant tout des machines allumées. Notre algorithme \sel{} en tire les
conséquences : il choisit exactement combien de VM de chaque type allumer, puis
répartit les tâches avec Max-Min.

En conditions réelles, il réduit le coût de 71\,\% et 29\,\% sur les petits lots et
le temps de 24\,\% sur le gros lot, par rapport à Round Robin. Le modèle prédit les
temps à environ $\pm 10\,\%$ quand les VM fonctionnent normalement. La principale
limite est la variabilité des VM à cœur partagé, qui peut ralentir une machine de
moitié : c'est elle qu'il faudrait traiter en priorité, par un rééquilibrage
dynamique ou une mesure des VM juste avant de planifier.


\appendix
% ============================================================================
\section{Reproduire les résultats}
\label{sec:repro}
% ============================================================================

\lstset{columns=flexible, keepspaces=true}

\paragraph{Compilation et tests (sans GCP).}
\begin{lstlisting}
make                # workloads locaux + orchestrator, scheduler, gentasks, gantt
make linux          # workloads pour les VM (linux/amd64)
make test           # tests du modèle et des algorithmes
\end{lstlisting}

\paragraph{Simulation (sans GCP).}
\begin{lstlisting}
./bin/gentasks -n 500 -seed 42
./bin/scheduler -compare -tasks data/tasks_500.csv
./bin/scheduler -explain -tasks data/tasks_500.csv
./bin/scheduler -compare -tasks data/tasks.csv     # + optimum exact
./bin/scheduler -algo select-maxmin -tasks data/tasks_500.csv \
                -out data/alloc.csv
\end{lstlisting}

\paragraph{Mesures et expériences sur GCP.}
\begin{lstlisting}
./scripts/manage_vms.sh start
./scripts/deploy.sh
./scripts/bench_vms.sh       # benchmark -> data/vm_performance.csv
./scripts/exp_workers.sh     # parallélisme réel de chaque VM
./scripts/exp_real.sh        # RR contre select-maxmin, 3 répétitions
./bin/orchestrator -alloc data/alloc.csv -out results/run.csv
./scripts/manage_vms.sh stop # toujours à la fin
\end{lstlisting}

\paragraph{Analyse des exécutions réelles (sans GCP).}
\begin{lstlisting}
./scripts/precision.sh results/exp_real/*/   # makespan prédit / mesuré
./bin/gantt results/exp_real_all/tasks/rr_r1.csv \
            results/exp_real_all/tasks/select-maxmin_r1.csv
\end{lstlisting}

\section{Organisation du dépôt}

\begin{lstlisting}
workloads/<nom>/main.go  les 4 workloads
orchestrator/            l'orchestrateur
scheduler/model/         modèle : prédiction, simulation, coût, score
scheduler/algorithms/    RR, MET, MCT, Min-Min, Max-Min, select-maxmin, optimal
scheduler/cmd/           programmes scheduler, gentasks et gantt
scripts/                 VM, déploiement, benchmark, expériences
data/                    tâches, machines, mesures, allocations
results/                 benchmarks, simulations, exécutions réelles
docs/RESULTATS.md        journal détaillé de toutes les mesures
\end{lstlisting}


\end{document}
